% Generated draft: requires mathematical and source review.
\documentclass[11pt]{article}
\usepackage[margin=25mm]{geometry}
\usepackage{fontspec}
\setmainfont{Latin Modern Roman}
\usepackage{amsmath,amssymb,mathtools,unicode-math}
\setmathfont{Latin Modern Math}
\usepackage{enumitem,needspace,xurl,hyperref}
\setlength{\parindent}{0pt}
\setlength{\parskip}{5pt}
\setlength{\emergencystretch}{3em}
\newcommand{\N}{\mathbb{N}}
\newcommand{\Z}{\mathbb{Z}}
\newcommand{\Q}{\mathbb{Q}}
\newcommand{\R}{\mathbb{R}}
\newcommand{\C}{\mathbb{C}}
\newcommand{\F}{\mathbb{F}}
\newcommand{\A}{\mathbb{A}}
\newcommand{\PP}{\mathbb{P}}
\newcommand{\E}{\mathbb{E}}
\newcommand{\eps}{\varepsilon}
\newcommand{\dd}{\,\mathrm{d}}
\begin{document}
\section*{UnsolvedMath ID: 30005354}
\textbf{Source problem code:} catalogue entry
\subsection{Definitions and mathematical statement}
Let \(k\) be an integer with \(k\ge 2\). Working in the usual finite simple undirected graph setting, let \(G=(V,E)\), let \(n=|V|>0\), let \(\delta(G)=\min_{v\in V}\deg_G(v)\), and let \(\chi(G)\) be the least number of colors in a proper vertex coloring. Write \(C_\ell\) for the cycle of length \(\ell\). The forbidden-cycle condition is
\[
C_{2j+1}\not\subseteq G\qquad(1\le j\le k-1),
\]
where containment means containment as a subgraph, not necessarily as an induced subgraph. Equivalently, every odd cycle in \(G\), if any exists, has length at least \(2k+1\).

The question attributed to Julia B\"ottcher asks whether these graphs can have \(\chi(G)\ge4\) and minimum degree at least a positive constant times \(n/k\), with the constant independent of \(k\) and \(n\). An explicit asymptotic interpretation of this target is:
\[
\begin{gathered}
\exists c>0\;\exists k_0\in\Z,\ k_0\ge2\;\forall k\in\Z,\ k\ge k_0\;\exists G\text{ finite simple undirected},\\
n=|V(G)|>0,\qquad
\chi(G)\ge4,\qquad
\delta(G)\ge\frac{c}{k}n,\qquad
C_{2j+1}\not\subseteq G\quad(1\le j\le k-1).
\end{gathered}
\]
This display makes explicit a uniform asymptotic lower bound. The supplied wording does not specify whether existence is required for some \(n\), arbitrarily large \(n\), or every sufficiently large \(n\), nor does it specify the initial range of \(k\). Thus the display records a minimal asymptotic interpretation, and its quantifier over graph order requires review against the original report.

The original excerpt asks for \(f(k)\) ``at least linear in \(1/k\)'', interpreted here as \(f(k)\ge c/k\) for all sufficiently large \(k\). The cleaned statement instead says ``of order \(1/k\)'', which ordinarily means that constants \(c,C>0\) satisfy \(c/k\le f(k)\le C/k\) for all sufficiently large \(k\). Because \(f(k)\) occurs as a lower degree bound, the existence target can be expressed using \(f(k)=c/k\); the wording difference is nevertheless retained for review.

The supplied original excerpt reports existence for
\[
f(k)=\frac{3}{2k^2+k+1}
\]
and reports nonexistence for
\[
f(k)=\frac{1}{2k}\qquad\text{when }k>7000,
\]
with an unresolved reference labeled [1]. These are supplied background claims, not independently verified results in this draft; their precise graph-order quantifiers are not given in the excerpt.

The original field also contains an apparently unrelated sentence about forcing an OCR-rendered difference resembling \(\alpha_{\pi_1}(G)-\alpha_{\mathbb F_2}(G)\) to tend to infinity using a bounded-degree triangulation of \(\mathbb{RP}^{2}\) and disjoint unions. Neither invariant is defined in the supplied record, and its relationship to the graph question is unspecified. It cannot be incorporated as a precise hypothesis or conclusion.
\subsection{Short English statement}
The question asks whether graphs with no odd cycle shorter than \(2k+1\) can require at least four colors while every vertex has degree at least \(cn/k\), for one positive constant \(c\) independent of \(k\). The supplied record labels the problem open. The graph-order quantifier, the initial range of \(k\), and a seemingly unrelated appended sentence require human review.
\subsection{Sources}
\begin{itemize}
\item Supplied local UnsolvedMath registry record, permanent ID \(30005354\), title: High-Chromatic Odd-Cycle-Free Graphs with Large Minimum Degree; exact source code \texttt{catalogue entry}. The original excerpt attributes Question 13 to Julia B\"ottcher; the supplied proposed year is 2023.
\item Supplied report citation: \emph{Combinatorics} (2023), Oberwolfach Reports, DOI 10.4171/owr/2023/1, pp. 5--89. \url{https://doi.org/10.4171/owr/2023/1}. This is also the supplied literature source labeled Oberwolfach Reports 2023/1.
\item Supplied literature source: Bottcher--Han--Kohayakawa--Person, \emph{Odd-Cycle-Free Graphs}. \url{https://arxiv.org/abs/2302.01875}.
\item Supplied literature source: Yuan--Peng, \emph{European Journal of Combinatorics}. \url{https://doi.org/10.1016/j.ejc.2025.104143}.
\item Literature checked at: \texttt{2026-08-22}. This date and the following assessment are supplied dated curation, not independently verified: the record labels the problem open, describes known constructions as having degree proportion of order \(1/k^2\), and describes minimum-degree results as forcing 3-colorability above roughly \(1/(2k)\), while leaving the constant-over-\(k\) existence question unresolved. The background explicitly requests independent verification of time-sensitive status and cited claims.
\item The registry labels the statement \texttt{corrected\_verified} and says it was checked against the official report. This is supplied verification metadata, not a verification performed for this draft. The original excerpt's reference [1] is not bibliographically identified in the supplied record.
\end{itemize}
\end{document}
